\documentclass[conference]{IEEEtran}

\usepackage{graphicx}
\usepackage{subfigure}
\usepackage{amsmath} % assumes amsmath package installed
\usepackage{amssymb} 
\usepackage{array}
\usepackage{tikz,circuitikz}
\usetikzlibrary{fit}
\usepackage{blindtext}
\usepackage{color}
\usepackage{siunitx}
\usepackage{float} % enables [H] float placement

\def\BibTeX{{\rm B\kern-.05em{\sc i\kern-.025em b}\kern-.08em
    T\kern-.1667em\lower.7ex\hbox{E}\kern-.125emX}}


\usepackage{mwe}
\usepackage{fancyhdr}
\fancypagestyle{firststyle}
{
	\fancyhf[C]{\fontsize{8}{10} \selectfont \textit{} }
	\fancyfoot[C]{}
}


\hyphenation{op-tical net-works semi-conduc-tor}


\begin{document}
%
% paper title
% Titles are only capitalized in the first letter.
% Linebreaks \\ can be used within to get better formatting as desired.
% Do not put math or special symbols in the title.
\title{Analsis de un rectifacdor de media onda con carga RL}

\author{
\IEEEauthorblockN{Yosniel Agüero}
\IEEEauthorblockA{\textit{Universidad de Guadalajara}\\
MCIE\\
Guadalajara, México\\
\scriptsize{\textit{yosniel.aguero9368@alumnos.udg.mx}}}
\and
\IEEEauthorblockN{Glader Hernandez}
\IEEEauthorblockA{\textit{Universidad de Guadalajara}\\
MCIE\\
Guadalajara, México\\
\scriptsize{\textit{glader.hernandez9367@alumnos.udg.mx}}}
\and
\IEEEauthorblockN{Gary Sosa}
\IEEEauthorblockA{\textit{Universidad de Guadalajara}\\
MCIE\\
Guadalajara, México\\
\scriptsize{\textit{gary.sosa9369@alumnos.udg.mx}}}
\and
\IEEEauthorblockN{Ulrik Wong}
\IEEEauthorblockA{\textit{Universidad de Guadalajara}\\
MCIE\\
Guadalajara, México\\
\scriptsize{\textit{ulrik.wong7998@alumnos.udg.mx}}}
}


\maketitle

\thispagestyle{firststyle}
\renewcommand{\headrulewidth}{0in}
\pagestyle{empty}


\pagestyle{fancy}
\chead{\fontsize{8}{10} \selectfont \textit{} }
\pagenumbering{gobble}



% As a general rule, do not put math, special symbols or citations
% in the abstract
\begin{abstract}
This document is a model and instructions for \LaTeX.
This and the IEEEtran.cls file define the components of your document [title, text, heads, etc.]. *CRITICAL: Do Not Use Symbols, Special Characters, Footnotes, 
or Math in Document title or Abstract.
\end{abstract}


\IEEEpeerreviewmaketitle



\section{Introduction}
This document is a model and instructions for \LaTeX.
Please observe the report page limits.
 
\begin{figure}[ht]
	\centering
	\begin{circuitikz}[american, cute inductors, american currents, american voltages]
		\ctikzset{diodes/scale=0.5, sources/scale=0.75, resistors/scale=0.5, inductors/scale=0.75}
		\draw (0,0) to[sinusoidal voltage source, v^<=$v_i$] (0,2) to [full diode, l_=$D_1$, v^=$v_{D}$, i_=$i_D$](2,2);
		\draw (2,0) to[full diode, l^=$D_2$, v_<=$v_{oi}$](2,2) to[L, l_=$L$, i_=$i_L$, v^=$v_L$] (4,2);
		\draw (4,2)-- (5,2) to[R, l_=$R$, v^=$v_R$, i>_=$i_R$] (5,0);
		\draw (0,0) -- (5,0);
	\end{circuitikz}
\end{figure}
\section{Análisis del rectificador monofásico de media onda con carga \(R\!-\!L\) y diodo de corrida libre}
En lo que sigue consideramos la fuente sinusoidal
\[
v_i(t)=V_m\sin(\omega t),
\]
una impedancia serie \(R\) y \(L\) (carga), y dos diodos ideales:
\begin{itemize}
  \item \(D_1\): rectificador que conecta la fuente a la carga (ánodo en la fuente, cátodo en el nodo de carga).
  \item \(D_2\): diodo de \emph{flyback } en paralelo a la carga (orientación tal que permite mantener la corriente inductiva cuando la fuente deja de alimentar).
\end{itemize}

\subsection{Estados asumidos de los diodos}

Definición de los 4 estados considerados y denotamos ON como  conducción (diodo polarizado directamente) y OFF como bloqueo (polarizado inversamente).
\[
\begin{array}{ll}
\text{Estado A: } & (D_1=\mathrm{ON},\; D_2=\mathrm{OFF}) \\
\text{Estado B: } & (D_1=\mathrm{OFF},\; D_2=\mathrm{ON}) \\
\text{Estado C: } & (D_1=\mathrm{OFF},\; D_2=\mathrm{OFF}) \\
\text{Estado D: } & (D_1=\mathrm{ON},\; D_2=\mathrm{ON}) \\
\end{array}
\]


\begin{itemize}
  \item \textbf{Operación de cada estado asumido :}
    \begin{enumerate}
      \item Estado A: Ocurre cuando la tensión instantánea de la fuente tiende a \emph{polarizar positivamente} \(D_1\) y puede imponer una tensión mayor en el nodo de carga que la necesaria para forzar \(D_2\) en conducción inversa. \(D_1\) conduce si \(\;v_i(t) > v_R(t)\) (ánodo de \(D_1\) más positivo que su cátodo). En la práctica, con diodos ideales y caída nula, el encendido ocurre cuando \(v_i(t)\) supera la tensión instantánea necesaria para mantener la corriente \(i(t)>0\).
      \item Estado B: Ocurre cuando la fuente no sostiene la corriente inductiva, pero la inercia del inductor mantiene corriente positiva; entonces \(D_2\) ofrece el camino de libre. \(D_2\) conduce si la polaridad en la carga, hace que el ánodo de \(D_2\) sea más positivo que su cátodo, es decir, cuando la inercia del inductor empuja la corriente y la tensión en la carga favorece la conducción por \(D_2\). Con la orientación antiparalela típica, esto sucede cuando \(v_i(t)\) cae por debajo de \(v_o(t)\) y la corriente inductiva es positiva.
      \item Estado C: Estado de no-conducción es válido si la corriente ha decaído a cero (\(i(t)=0\)  ) y ambas tensiones en los ánodos/cátodos de los diodos dejan a ambos inversamente polarizados.
      \item Estado D: Si los diodos son ideales y están orientados en la configuración habitual (antiparalelo para \(D_2\)), la conducción simultánea tiende a producir una contradicción en las polaridades o a generar un lazo de circulación que, en el mejor de los casos, fuerza igualar tensiones entre nodos y, en el peor, crea una trayectoria de cortocircuito entre la fuente y el retorno. Por tanto en la práctica se considera no-sostenible como estado permanente y debe justificarse con cuidado.
\end{enumerate}
\end{itemize}

 


\subsection{Modelo en espacio de estados de los estados válidos }
En este sistema la única variable de estado es la corriente \(i(t)\) de la carga \(R\!-\!L\). Usaremos la variable de estado:
\[
x(t)=i(t).
\]
A continuación planteamos las ecuaciones de estado para cada configuración válida (y para la configuración asumida D).

\paragraph{Estado A: \(D_1=\mathrm{ON},\;D_2=\mathrm{OFF}\) (circuito L–R  con fuente de alimentación).}
La ecuación diferencial y la representación en espacio de estados:
\[
L\frac{di}{dt} + R\,i(t) = v_s(t) = V_m\sin(\omega t).
\]
\[
\begin{bmatrix}\dot{x}(t)\end{bmatrix}
=
\begin{bmatrix}-\dfrac{R}{L}\end{bmatrix}
\begin{bmatrix}x(t)\end{bmatrix}
+
\begin{bmatrix}\dfrac{1}{L}\end{bmatrix}
\qquad
y(t)=\begin{bmatrix}R\end{bmatrix}\begin{bmatrix}x(t)\end{bmatrix}.
\]
Solución general (en \(t_0\) con condición \(i(t_0)=I_0\)):
\[
i(t)=e^{-\frac{R}{L}(t-t_0)}I_0 + \frac{V_m}{L}\int_{t_0}^{t} e^{-\frac{R}{L}(t-\tau)}\sin(\omega \tau)\,d\tau.
\]
La parte forzada en régimen permanente (cuando el tiempo es mucho mayor que \(\tfrac{L}{R}\)) tiene la forma:
\[
i_p(t)=\frac{V_m}{\sqrt{R^2+(\omega L)^2}}\sin(\omega t - \phi),\qquad \phi=\arctan\!\frac{\omega L}{R}.
\]

\paragraph{Estado B: (diodo \emph{flyback } — circuito L–R  sin fuente).}
Ecuación diferencial:
\[
L\frac{di}{dt} + R\,i(t) = 0
\]
\[
\begin{bmatrix}\dot{x}(t)\end{bmatrix}
=
\begin{bmatrix}-\dfrac{R}{L}\end{bmatrix}
\begin{bmatrix}x(t)\end{bmatrix}
\qquad
y(t)=\begin{bmatrix}R\end{bmatrix}\begin{bmatrix}x(t)\end{bmatrix}.
\]
La solucion general es:
\[
i(t)=I_1\,e^{-\frac{R}{L}(t-t_1)},
\]
donde \(I_1\) es la corriente en el instante \(t_1\) en que comienza el freewheeling.
\paragraph{Estado C: (no-conducción).}
Si \(x(t)=i(t)=0\) este estado es consistente; la dinámica es \(i\equiv 0\). Si el inductor intenta imponer \(i\neq0\). 


\section{Ease of Use}

\subsection{Maintaining the Integrity of the Specifications}

The IEEEtran class file is used to format your document and style the text. All margins, 
column widths, line spaces, and text fonts are prescribed; please do not 
alter them. You may note peculiarities. For example, the head margin
measures proportionately more than is customary. This measurement 
and others are deliberate, please do not revise any of the current designations.

\subsection{Subsection Heading Here}
Subsection text here.

\section{Prepare Your Document Before Styling}
Before you begin to format your document, first write and save the content as a 
separate text file. Complete all content and organizational editing before 
formatting. Please note sections \ref{AA}--\ref{SCM} below for more information on 
proofreading, spelling and grammar.

Keep your text and graphic files separate until after the text has been 
formatted and styled. Do not number text heads---{\LaTeX} will do that 
for you.

\subsection{Abbreviations and Acronyms}\label{AA}
Define abbreviations and acronyms the first time they are used in the text, 
even after they have been defined in the abstract. Abbreviations such as 
IEEE, SI, MKS, CGS, ac, dc, and rms do not have to be defined. Do not use 
abbreviations in the title or heads unless they are unavoidable.

\subsection{Units}
\begin{itemize}
\item Use either SI (MKS) or CGS as primary units. (SI units are encouraged.) English units may be used as secondary units (in parentheses). An exception would be the use of English units as identifiers in trade, such as ``3.5-inch disk drive''.
\item Avoid combining SI and CGS units, such as current in amperes and magnetic field in oersteds. This often leads to confusion because equations do not balance dimensionally. If you must use mixed units, clearly state the units for each quantity that you use in an equation.
\item Do not mix complete spellings and abbreviations of units: ``Wb/m\textsuperscript{2}'' or ``webers per square meter'', not ``webers/m\textsuperscript{2}''. Spell out units when they appear in text: ``. . . a few henries'', not ``. . . a few H''.
\item Use a zero before decimal points: ``0.25'', not ``.25''. Use ``cm\textsuperscript{3}'', not ``cc''.)
\end{itemize}

\subsection{Equations}
Number equations consecutively. To make your 
equations more compact, you may use the solidus (~/~), the exp function, or 
appropriate exponents. Italicize Roman symbols for quantities and variables, 
but not Greek symbols. Use a long dash rather than a hyphen for a minus 
sign. Punctuate equations with commas or periods when they are part of a 
sentence, as in:
\begin{equation}
a+b=\gamma.\label{eq}
\end{equation}

Be sure that the 
symbols in your equation have been defined before or immediately following 
the equation. Use ``\eqref{eq}'', not ``Eq.~\eqref{eq}'' or ``equation \eqref{eq}'', except at 
the beginning of a sentence: ``Equation \eqref{eq} is . . .''

\subsection{\LaTeX-Specific Advice}

Please use ``soft'' (e.g., \verb|\eqref{Eq}|) cross references instead
of ``hard'' references (e.g., \verb|(1)|). That will make it possible
to combine sections, add equations, or change the order of figures or
citations without having to go through the file line by line.

Please don't use the \verb|{eqnarray}| equation environment. Use
\verb|{align}| or \verb|{IEEEeqnarray}| instead. The \verb|{eqnarray}|
environment leaves unsightly spaces around relation symbols.

Please note that the \verb|{subequations}| environment in {\LaTeX}
will increment the main equation counter even when there are no
equation numbers displayed. If you forget that, you might write a
document in which the equation numbers skip from (17) to (20), causing
the reader to wonder if you've discovered a new method of
counting.

{\BibTeX} does not work by magic. It doesn't get the bibliographic
data from thin air but from .bib files. If you use {\BibTeX} to produce a
bibliography you must include the .bib files. 

{\LaTeX} can't read your mind. If you assign the same label to a
subsubsection and a table, you might find that Table I has been cross
referenced as Table IV-B3. 

{\LaTeX} does not have precognitive abilities. If you put a
\verb|\label| command before the command that updates the counter it's
supposed to be using, the label will pick up the last counter to be
cross referenced instead. In particular, a \verb|\label| command
should not go before the caption of a figure or a table.

Do not use \verb|\nonumber| inside the \verb|{array}| environment. It
will not stop equation numbers inside \verb|{array}| (there won't be
any anyway) and it might stop a wanted equation number in the
surrounding equation.

\subsection{Some Common Mistakes}\label{SCM}
\begin{itemize}
\item The word ``data'' is plural, not singular.
\item The subscript for the permeability of vacuum $\mu_{0}$, and other common scientific constants, is zero with subscript formatting, not a lowercase letter ``o''.
\item In American English, commas, semicolons, periods, question and exclamation marks are located within quotation marks only when a complete thought or name is cited, such as a title or full quotation. When quotation marks are used, instead of a bold or italic typeface, to highlight a word or phrase, punctuation should appear outside of the quotation marks. A parenthetical phrase or statement at the end of a sentence is punctuated outside of the closing parenthesis (like this). (A parenthetical sentence is punctuated within the parentheses.)
\item A graph within a graph is an ``inset'', not an ``insert''. The word alternatively is preferred to the word ``alternately'' (unless you really mean something that alternates).
\item Do not use the word ``essentially'' to mean ``approximately'' or ``effectively''.
\item In your document title, if the words ``that uses'' can accurately replace the word ``using'', capitalize the ``u''; if not, keep using lower-cased.
\item Be aware of the different meanings of the homophones ``affect'' and ``effect'', ``complement'' and ``compliment'', ``discreet'' and ``discrete'', ``principal'' and ``principle''.
\item Do not confuse ``imply'' and ``infer''.
\item The prefix ``non'' is not a word; it should be joined to the word it modifies, usually without a hyphen.
\item There is no period after the ``et'' in the Latin abbreviation ``et al.''.
\item The abbreviation ``i.e.'' means ``that is'', and the abbreviation ``e.g.'' means ``for example''.
\end{itemize}
An excellent style manual for science writers is \cite{b7}.

\subsection{Authors and Affiliations}
\textbf{The class file is designed for, but not limited to, six authors.} A 
minimum of one author is required. Author names 
should be listed starting from left to right and then moving down to the 
next line. Names should not be listed in columns nor group by 
affiliation. Please keep your affiliations as succinct as possible (for 
example, do not differentiate among departments of the same organization).

\subsection{Identify the Headings}
Headings, or heads, are organizational devices that guide the reader through 
your document. There are two types: component heads and text heads.

Component heads identify the different components of your document and are not 
topically subordinate to each other. Examples include Acknowledgments and 
References and, for these, the correct style to use is ``Heading 5''. Use 
``figure caption'' for your Figure captions, and ``table head'' for your 
table title. Run-in heads, such as ``Abstract'', will require you to apply a 
style (in this case, italic) in addition to the style provided by the drop 
down menu to differentiate the head from the text.

Text heads organize the topics on a relational, hierarchical basis. For 
example, the document title is the primary text head because all subsequent 
material relates and elaborates on this one topic. If there are two or more 
sub-topics, the next level head (uppercase Roman numerals) should be used 
and, conversely, if there are not at least two sub-topics, then no subheads 
should be introduced.

\subsection{Figures and Tables}
\paragraph{Positioning Figures and Tables} Place figures and tables at the top and 
bottom of columns. Avoid placing them in the middle of columns. Large 
figures and tables may span across both columns. Figure captions should be 
below the figures; table heads should appear above the tables. Insert 
figures and tables after they are cited in the text. Use the abbreviation 
``Fig.~\ref{fig}'', even at the beginning of a sentence.

\begin{table}[htbp]
\caption{Table Type Styles}
\begin{center}
\begin{tabular}{|c|c|c|c|}
\hline
\textbf{Table}&\multicolumn{3}{|c|}{\textbf{Table Column Head}} \\
\cline{2-4} 
\textbf{Head} & \textbf{\textit{Table column subhead}}& \textbf{\textit{Subhead}}& \textbf{\textit{Subhead}} \\
\hline
copy& More table copy$^{\mathrm{a}}$& &  \\
\hline
\multicolumn{4}{l}{$^{\mathrm{a}}$Sample of a Table footnote.}
\end{tabular}
\label{tab1}
\end{center}
\end{table}

\begin{figure}[htbp]
\centerline{\includegraphics{fig1}}
\caption{Example of a figure caption.}
\label{fig}
\end{figure}

Figure Labels: Use 8 point Times New Roman for Figure labels. Use words 
rather than symbols or abbreviations when writing Figure axis labels to 
avoid confusing the reader. As an example, write the quantity 
``Magnetization'', or ``Magnetization, M'', not just ``M''. If including 
units in the label, present them within parentheses. Do not label axes only 
with units. In the example, write ``Magnetization (A/m)'' or ``Magnetization 
\{A[m(1)]\}'', not just ``A/m''. Do not label axes with a ratio of 
quantities and units. For example, write ``Temperature (K)'', not 
``Temperature/K''.

\section*{References}

Please number citations consecutively within brackets \cite{b1}. The 
sentence punctuation follows the bracket \cite{b2}. Refer simply to the reference 
number, as in \cite{b3}---do not use ``Ref. \cite{b3}'' or ``reference \cite{b3}'' except at 
the beginning of a sentence: ``Reference \cite{b3} was the first $\ldots$''

Number footnotes separately in superscripts. Place the actual footnote at 
the bottom of the column in which it was cited. Do not put footnotes in the 
abstract or reference list. Use letters for table footnotes.

Unless there are six authors or more give all authors' names; do not use 
``et al.''. Papers that have not been published, even if they have been 
submitted for publication, should be cited as ``unpublished'' \cite{b4}. Papers 
that have been accepted for publication should be cited as ``in press'' \cite{b5}. 
Capitalize only the first word in a paper title, except for proper nouns and 
element symbols.

For papers published in translation journals, please give the English 
citation first, followed by the original foreign-language citation \cite{b6}.

\section{Conclusions}
The conclusion goes here.

\begin{thebibliography}{plain}
\bibitem{b1} G. Eason, B. Noble, and I. N. Sneddon, ``On certain integrals of Lipschitz-Hankel type involving products of Bessel functions,'' Phil. Trans. Roy. Soc. London, vol. A247, pp. 529--551, April 1955.
\bibitem{b2} J. Clerk Maxwell, A Treatise on Electricity and Magnetism, 3rd ed., vol. 2. Oxford: Clarendon, 1892, pp.68--73.
\bibitem{b3} I. S. Jacobs and C. P. Bean, ``Fine particles, thin films and exchange anisotropy,'' in Magnetism, vol. III, G. T. Rado and H. Suhl, Eds. New York: Academic, 1963, pp. 271--350.
\bibitem{b4} K. Elissa, ``Title of paper if known,'' unpublished.
\bibitem{b5} R. Nicole, ``Title of paper with only first word capitalized,'' J. Name Stand. Abbrev., in press.
\bibitem{b6} Y. Yorozu, M. Hirano, K. Oka, and Y. Tagawa, ``Electron spectroscopy studies on magneto-optical media and plastic substrate interf	ace,'' IEEE Transl. J. Magn. Japan, vol. 2, pp. 740--741, August 1987 [Digests 9th Annual Conf. Magnetics Japan, p. 301, 1982].
\bibitem{b7} M. Young, The Technical Writer's Handbook. Mill Valley, CA: University Science, 1989.
\end{thebibliography}

\vspace{12pt}
\color{red}
This report template contains guidance text for composing and formatting documents. Please ensure that all template text is removed from your document prior to submission of your work. Failure to remove the template text from your document will be considered as a lack of editing effort.

\end{document}